% !PRINTED 218
\weberpage{218}

$\Theta$ ist eine sogenannte zerlegbare Form $\frac{1}{2}n(n-1)^{\mathrm{ten}}$ Grades in den Variablen $t$; denn sie lässt sich nach (7) in lauter lineare Factoren zerlegen, die freilich nicht rational in den $a$ sind.

Wir haben schon im \S.\,52 darauf hingewiesen, dass bei der Tschirnhausen-Transformation auch $x$ rational durch $y$ ausdrückbar ist; und dasselbe gilt also auch für jede rationale Function von $x$.

Betrachten wir irgend eine solche Function $\varphi(x)$, die auch noch die Coëfficienten $a$ und $t$ enthalten kann, aber immer für beide Arten von Variablen ganz und homogen vorausgesetzt sei, so können wir setzen

\weq{9}{\varphi(x) = C_0 + C_1y + C_2y^2 + \cdots + C_{n-1}y^{n-1},}

worin die $C_0, C_1, \ldots C_{n-1}$ rational in $a$ und in $t$ ausdrückbar sind.

Stellen wir die Gleichung (9) für $x=x_1,x_2\ldots x_n$, $y=y_1,y_2,\ldots y_n$ auf, so erhalten wir für die Bestimmung der $C$ ein System linearer Gleichungen, dessen Determinante

\weq{10}{\begin{vmatrix} 1, & y_1, & y_1^2 & \ldots & y_1^{n-1} \\ 1, & y_2, & y_2^2 & \ldots & y_2^{n-1} \\ \cdot & \cdot & \cdot & \cdot & \cdot \\ 1, & y_n, & y_n^2 & \ldots & y_n^{n-1} \end{vmatrix},}

also gleich $\Theta\sqrt{D}$ ist. Der Zähler des Ausdruckes für $C_\nu$ geht aus der Determinante (10) hervor, indem man die Elemente der $(\nu+1)^{\mathrm{ten}}$ Colonne durch $\varphi(x_1), \varphi(x_2)\ldots\varphi(x_n)$ ersetzt, und folglich hat der Zähler, wenn man Alles durch die $x_i$ ausdrückt, den Factor $\sqrt{D}$ (weil er verschwindet, wenn zwei $x_i$ einander gleich werden).

Wir können also setzen

\[ \Theta C_\nu = Q_\nu, \]

worin $Q_\nu$ eine ganze rationale Function der $t$ und der $a$ ist, die höchstens noch eine Potenz von $a_0$ im Nenner enthalten kann. Wir bekommen dann

\weq{11}{\Theta\varphi(x) = Q_0 + Q_1y + Q_2y^2 + \cdots + Q_{n-1}y^{n-1}.}

Die Coëfficienten $Q_\nu$ sind aber keine Invarianten und ihre Berechnung ist in den meisten Fällen schwierig.
